% =====================================================================
%  CHƯƠNG 5. KIỂM THỬ VÀ KẾT QUẢ THỰC NGHIỆM
% =====================================================================
\chapter{KIỂM THỬ VÀ KẾT QUẢ THỰC NGHIỆM}
\label{chap:5}

% ---------------------------------------------------------------------
\section{Bộ 10 ca kiểm thử tự động}
\label{sec:5.1}

\begin{table}[H]
\centering
\caption{Bộ 10 ca kiểm thử tự động}
\label{tab:10ca}
\begin{tabular}{|c|p{5.2cm}|p{3.4cm}|p{3.4cm}|c|}
\hline
\textbf{STT} & \textbf{Nội dung kiểm thử} & \textbf{Đầu vào} &
\textbf{Kết quả mong đợi} & \textbf{Đạt}\\ \hline
1  & & & & \\ \hline
2  & & & & \\ \hline
3  & & & & \\ \hline
4  & & & & \\ \hline
5  & & & & \\ \hline
6  & & & & \\ \hline
7  & & & & \\ \hline
8  & & & & \\ \hline
9  & & & & \\ \hline
10 & & & & \\ \hline
\end{tabular}
\end{table}

\canviet{Điền theo \code{test\_hamming.py}. Gợi ý nhóm ca cần có: mã hóa đúng
16 từ mã; $G\HT = 0$; không lỗi; sửa đúng mọi lỗi đơn (112/112); hai lỗi ở
(7,4) bị sửa nhầm (336/336); hai lỗi ở (8,4) được phát hiện; chuỗi ký tự đầu --
cuối; chuỗi rỗng; ký tự UTF-8; đầu vào sai kiểu/độ dài. Kèm ảnh hoặc log kết
quả chạy.}

% ---------------------------------------------------------------------
\section{Phân loại kết quả theo số lỗi trong khối}
\label{sec:5.2}

Bảng~\ref{tab:theosoloi} là kết quả vét cạn: với mỗi dữ liệu $u$ (16 giá trị)
và mỗi vectơ lỗi $e \in \GFn{7}$ (128 giá trị), giải mã $r = uG \xor e$ và so
sánh với $u$.

\begin{table}[H]
\centering
\caption{Kết quả giải mã (7,4) theo số bit lỗi $k$ trong một khối (vét cạn $16 \times 128$ mẫu)}
\label{tab:theosoloi}
\begin{tabular}{|c|c|c|c|}
\hline
\textbf{Số lỗi $k$} & \textbf{Số mẫu} & \textbf{Hội chứng $= 0$} &
\textbf{Giải mã đúng}\\ \hline
0 & 16  & 16  & 16\\
1 & 112 & 0   & 112\\
2 & 336 & 0   & 0\\
3 & 560 & 112 & 0\\
4 & 560 & 112 & 0\\
5 & 336 & 0   & 0\\
6 & 112 & 0   & 0\\
7 & 16  & 16  & 0\\ \hline
\textbf{Tổng} & 2048 & 256 & 128\\ \hline
\end{tabular}
\end{table}

\canviet{Xác nhận lại bảng bằng chính chương trình của nhóm (số liệu trên đã
chạy vét cạn với bộ giải mã theo Bảng~\ref{tab:hoichung}). Nhận xét: giải mã
đúng khi và chỉ khi $k \le 1$; hội chứng bằng 0 với $k \ge 3$ xảy ra đúng khi
$e$ là một từ mã khác 0 (15 vectơ lỗi, nhân 16 dữ liệu).}

% ---------------------------------------------------------------------
\section{Mô phỏng chuỗi ký tự đầu -- cuối}
\label{sec:5.3}

\canviet{Chạy với một hoặc vài chuỗi cụ thể (ví dụ \code{"A"}, một câu tiếng
Việt có dấu) ở một số mức $p$: in chuỗi gửi, số bit, số khối bị lỗi, số khối
sửa đúng, chuỗi nhận được. Dán kết quả thật vào khối dưới.}

\begin{center}
\begin{minipage}{\linewidth}
\captionof{chuongtrinh}{Kết quả chạy mô phỏng chuỗi ký tự}\label{lst:ketqua53}
\end{minipage}
\end{center}
\begin{Verbatim}
# Dán nguyên văn đầu ra của chương trình tại đây
\end{Verbatim}

% ---------------------------------------------------------------------
\section{Tỉ lệ lỗi bit sau giải mã trên kênh BSC}
\label{sec:5.4}

\begin{table}[H]
\centering
\caption{Kết quả mô phỏng trên kênh BSC}
\label{tab:ber}
\begin{tabular}{|c|>{\centering\arraybackslash}p{3.4cm}|>{\centering\arraybackslash}p{3.8cm}|>{\centering\arraybackslash}p{3.8cm}|}
\hline
$p$ & \textbf{BER không mã hóa} & \textbf{BER sau giải mã (mô phỏng)} &
\textbf{Xác suất khối sai (lý thuyết)}\\ \hline
0{,}01 & & & 0{,}0020\\
0{,}05 & & & 0{,}0444\\
0{,}10 & & & 0{,}1497\\ \hline
\end{tabular}
\end{table}

Cột cuối tính theo \eqref{eq:pkhoi}:
$\Pr(\text{khối sai}) = 1 - (1-p)^7 - 7p(1-p)^6$.

\canviet{Điền BER mô phỏng (số khối, số lần lặp, \code{seed}); có thể thêm đồ
thị BER theo $p$ (thang log). Lưu ý: BER sau giải mã là tỉ lệ bit dữ liệu sai,
khác với xác suất khối sai ở cột cuối; chỉ dùng cột cuối để đối chiếu xu hướng.}

% ---------------------------------------------------------------------
\section{Nhận xét}
\label{sec:5.5}

\canviet{Đối chiếu kết quả thực nghiệm với lý thuyết Chương 1--2: sửa đúng mọi
lỗi đơn; hai lỗi bị sửa nhầm ở (7,4) nhưng được phát hiện ở (8,4); với $p$
nhỏ, mã hóa giảm đáng kể tỉ lệ lỗi; với $p$ lớn, lợi ích giảm.}
